\documentclass[12pt]{article}

\usepackage[a4paper, margin=2.5cm]{geometry}
\usepackage{amsmath}
\usepackage{amssymb}
\usepackage{amsthm}
\usepackage[colorlinks=true, linkcolor=blue, citecolor=blue, urlcolor=blue]{hyperref}
\usepackage{booktabs}
\usepackage{natbib}

\newtheorem{theorem}{Theorem}
\newtheorem{proposition}[theorem]{Proposition}
\newtheorem{corollary}[theorem]{Corollary}
\newtheorem{definition}[theorem]{Definition}
\newtheorem{remark}[theorem]{Remark}
\newtheorem{example}[theorem]{Example}
\newtheorem{lemma}[theorem]{Lemma}

\newcommand{\Ecal}{\mathcal{E}}
\newcommand{\Gcal}{\mathcal{G}}

\title{Dynamics and Field Equations\\from Event-Order Geometry}
\author{%
  Lee Smart\\[4pt]
  \textit{Institute of Vibrational Field Dynamics}\\[2pt]
  \texttt{contact@vibrationalfielddynamics.org}\\[2pt]
  \texttt{@vfd\_org}%
}
\date{April 2026}

\begin{document}

\maketitle

\noindent\textbf{Status:} Preprint (not peer-reviewed)

\begin{abstract}
Papers~XXIII--XXVI derived temporal ordering, observer kinematics, separation geometry, and three candidate graph-curvature indicators from event-order structure without assuming spacetime. The present paper takes the next step: introducing source candidates and exhibiting a discrete Poisson response on toy graphs.

A \textit{source field} $S(e)$ is defined from local event-content deviation, with three candidates (density, signed overlap, signed geodesic traffic) compared explicitly. An \textit{accessibility potential} $u(e)$ is introduced on the transition graph, governed by the discrete graph Laplacian equation $\Delta_\Gcal u(e) = S(e) - \bar{S}$ (with $u$ determined up to an additive constant). A potential-gradient curvature indicator $K_{\mathrm{pot}}$ is constructed from $u$; in the bowtie toy model below, $K_{\mathrm{pot}}$ is largest at the same event where Paper~XXVI's $K_{\mathrm{branch}}$ and $K_{\mathrm{geo}}$ are largest, with the caveat that the source used there and $K_{\mathrm{branch}}$ are built from the same degree data. No such comparison is claimed for the perturbation example. A global identity $\sum_e [S(e) - \bar{S}] = 0$ provides the compatibility condition required for the Poisson equation to be solvable.

Three toy models are worked explicitly: a uniform vacuum ($K_{5,5}$ pentagon truncation: zero source, constant potential, uniform betweenness), a static localised source (bowtie: nonzero potential, high $K_{\mathrm{pot}}$ at the bottleneck, and, independently, geodesic concentration at the bottleneck), and a perturbative model (adding one event to $K_{5,5}$: closed-form Poisson solve over four symmetry orbits, largest $K_{\mathrm{pot}}$ at the new event). In the bowtie example, large potential gradients coincide with concentrated geodesic traffic; in the perturbation example only new shortest paths are exhibited (no full betweenness recomputation is done). No general geodesic law from $u$ is derived.

The paper does not derive the Einstein equations or continuous field dynamics. It defines and analyzes a discrete graph-Poisson source-to-potential model on the event-order framework.
\end{abstract}

%==================================================================
\section{Introduction}\label{sec:intro}
%==================================================================

Paper~XXVI~\citep{SmartXXVI} constructed three candidate curvature indicators on the event-order framework: volume-growth distortion $K_{\mathrm{grow}}$ (reference-dependent), branching asymmetry $K_{\mathrm{branch}}$, and geodesic concentration $K_{\mathrm{geo}}$. In the flat $K_{5,5}$ toy model, all three vanish under the reference $V_{\mathrm{ref}} = \bar V$ at radius 1; in the asymmetric bowtie, $K_{\mathrm{branch}}$ and $K_{\mathrm{geo}}$ are nonzero (reference-independent) and $K_{\mathrm{grow}}$ is nonzero under the same reference choice at radius 1. But geometry without a source has no physics. The present paper addresses: \textit{what drives the distortion?}

\medskip
\noindent\textbf{Standing assumptions and imported notation.} Throughout this paper, $\Ecal$ is a finite event set, $\Gcal = (\Ecal, E_H)$ is the connected undirected Hasse graph of Paper~XXV~\citep{SmartXXV}, and $d_{\mathrm{trans}}$ is the associated shortest-path metric. Finiteness and connectedness of $\Gcal$ are standing; both toy models below are finite and connected by inspection. All sums over $\Ecal$ are finite.
\medskip


In general relativity, the answer is the stress-energy tensor: matter and energy tell geometry how to curve. In the event-order framework, the analogue is: \textit{non-uniform event content tells the accessibility geometry how to distort.}

\subsection*{Central principle}

\begin{quote}
Source is non-uniform event content. The field response is modelled as a geometric response of accessibility to that non-uniformity.
\end{quote}

\subsection*{Scope}

This paper constructs source candidates, poses a graph-Poisson field equation, proves the required compatibility and boundary-flux identities, and evaluates the construction on three toy graphs. It does not derive the Einstein equations, continuous dynamics, or stress-energy tensors. The graph-Poisson equation is introduced as a model on the event-order framework, not derived from prior VFD structure.

%==================================================================
\section{Source Quantities}\label{sec:source}
%==================================================================

Three candidate source constructions are developed and compared.

\subsection{Source A: Event density}

\begin{definition}[Density source]\label{def:rho}
Fix a radius $r \in \mathbb{Z}_{\geq 0}$ and a reference volume profile $V_{\mathrm{ref}}(r) > 0$ (e.g.\ the event-set average $\bar V(r)$). For each $e \in \Ecal$, define
\begin{equation}\label{eq:rho}
    S_\rho(e) \;=\; \frac{V_e(r)}{V_{\mathrm{ref}}(r)} - 1,
\end{equation}
where $V_e(r) = |B_r(e)|$ is the metric-ball size in the Hasse graph (Paper~XXVI). With $r$ and $V_{\mathrm{ref}}$ fixed once and for all, $S_\rho$ is a scalar field $\Ecal \to \mathbb{R}$ suitable for use as the right-hand side of the field equation below.
\end{definition}

$S_\rho = 0$: $e$'s radius-$r$ ball matches the reference (vacuum-like at that radius). $S_\rho > 0$: denser than reference. $S_\rho < 0$: sparser than reference.

\subsection{Source B: Overlap multiplicity}

\begin{definition}[Overlap source]\label{def:overlap_source}
For each $e \in \Ecal$, define
\begin{equation}
    S_\Omega(e) \;=\; \deg_\Gcal(e) - \overline{\deg},
\end{equation}
the signed degree deviation in $\Gcal$, where $\deg_\Gcal(e) = |\{e' \in \Ecal : d_{\mathrm{trans}}(e, e') = 1\}|$ and $\overline{\deg} = |\Ecal|^{-1} \sum_e \deg_\Gcal(e)$. By construction $\bar S_\Omega = 0$.
\end{definition}

This measures local interaction load: events with many admissible neighbors carry more ``content'' than average.

\subsection{Source C: Geodesic traffic}

\begin{definition}[Traffic source]\label{def:traffic}
\begin{equation}
    S_J(e) = \beta(e) - \bar{\beta}
\end{equation}
where $\beta(e)$ is the geodesic betweenness centrality from Paper~XXVI.
\end{definition}

Regions with concentrated geodesic traffic act as sources: they carry disproportionate flow, analogous to energy-momentum flux.

\begin{remark}[Source comparison]
$S_\rho$ measures static local density; $S_\Omega$ measures signed local connectivity excess; $S_J$ measures signed geodesic-traffic concentration. These are \emph{related to but not identical to} the curvature indicators of Paper~XXVI~\citep{SmartXXVI}: $K_{\mathrm{branch}}(e) = |\deg(e) - \overline{\deg}|$ is the \emph{absolute} degree deviation while $S_\Omega(e) = \deg(e) - \overline{\deg}$ is signed; $K_{\mathrm{geo}}(e) = (\beta(e) - \bar\beta)/\bar\beta$ is normalised while $S_J(e) = \beta(e) - \bar\beta$ is unnormalised. The sign information is essential for the Poisson equation below, which requires a signed right-hand side that integrates to zero. The framework permits any of these source choices in the field equation; in the worked examples we use whichever yields the cleanest closed-form, and we flag each choice explicitly.
\end{remark}

%==================================================================
\section{Accessibility Potential and Field Equation}\label{sec:field}
%==================================================================

\subsection{The graph Laplacian}

\begin{definition}[Graph Laplacian on $\Gcal$]\label{def:laplacian}
For the undirected transition graph $\Gcal = (\Ecal, E_H)$, the graph Laplacian acts on functions $f: \Ecal \to \mathbb{R}$ by:
\begin{equation}\label{eq:laplacian}
    (\Delta_\Gcal f)(e) = \sum_{e' \sim e} [f(e') - f(e)]
\end{equation}
where the sum is over neighbors of $e$ in $\Gcal$.
\end{definition}

\subsection{The discrete field equation}

\begin{definition}[Accessibility potential]\label{def:potential}
Given a source field $S: \Ecal \to \mathbb{R}$, the \textit{accessibility potential} is any function $u: \Ecal \to \mathbb{R}$ satisfying
\begin{equation}\label{eq:field}
    \boxed{\Delta_\Gcal\, u(e) = S(e) - \bar{S}}
\end{equation}
with $\bar{S} = |\Ecal|^{-1}\sum_{e'} S(e')$. The potential is determined only up to an additive constant (Proposition~\ref{prop:existence}); when a specific representative is required for arithmetic, we impose the zero-mean gauge $\sum_{e \in \Ecal} u(e) = 0$ or a pointwise gauge such as $u(e_0) = 0$ at a chosen base event $e_0$.
\end{definition}

\begin{remark}[Why $S - \bar{S}$]
The graph Laplacian satisfies $\sum_e (\Delta_\Gcal f)(e) = 0$ for any $f$. Therefore the equation $\Delta_\Gcal u = g$ has a solution only if $\sum g(e) = 0$. Subtracting $\bar{S}$ ensures this compatibility condition is satisfied.
\end{remark}

\begin{proposition}[Existence of the potential]\label{prop:existence}
Let $\Gcal$ be a finite connected undirected graph (in particular, $|\Ecal|$ is finite). The equation $\Delta_\Gcal u = S - \bar{S}$ has a solution $u: \Ecal \to \mathbb{R}$, unique up to an additive constant.
\end{proposition}

\begin{proof}
On a finite graph, the Laplacian $\Delta_\Gcal$ is a symmetric $|\Ecal| \times |\Ecal|$ matrix acting on $\mathbb{R}^\Ecal$, and its kernel equals $\mathrm{span}\{\mathbf{1}\}$ on a connected graph (the constant functions). The right-hand side satisfies $\sum_e [S(e)-\bar S] = 0$ (Proposition~\ref{prop:balance} below), so it lies in the image of $\Delta_\Gcal$ (which is the orthogonal complement of $\ker\Delta_\Gcal$). A solution therefore exists and is unique modulo constants.
\end{proof}

\subsection{Potential-gradient curvature}

The accessibility potential $u(e)$ is a scalar function on $\Ecal$; its local variation gives a natural candidate curvature indicator that sits alongside the three indicators of Paper~XXVI. We do not prove a general comparison theorem between $K_{\mathrm{pot}}$ (defined below) and $K_{\mathrm{grow}}, K_{\mathrm{branch}}, K_{\mathrm{geo}}$. In the bowtie toy model below, large $K_{\mathrm{pot}}$ occurs at the same event $b$ where Paper~XXVI found large $K_{\mathrm{branch}}$ and $K_{\mathrm{geo}}$; note that in that example $S = S_\Omega$ is built from signed degree deviation, so the $K_{\mathrm{pot}}$--$K_{\mathrm{branch}}$ alignment is a same-data alignment (both driven by $\deg - \overline{\deg}$), not an independent cross-check. The $K_{\mathrm{pot}}$--$K_{\mathrm{geo}}$ co-location is not a same-data alignment.

\begin{definition}[Potential-gradient curvature]\label{def:Kpot}
\begin{equation}
    K_{\mathrm{pot}}(e) = \sum_{e' \sim e} [u(e') - u(e)]^2
\end{equation}
the squared gradient of the accessibility potential at $e$.
\end{definition}

\begin{remark}
$K_{\mathrm{pot}} = 0$ iff $u$ is locally constant (no gradient, flat potential landscape). $K_{\mathrm{pot}} > 0$ indicates local potential variation: the source has created a ``hill'' or ``valley'' in accessibility that distorts the geometry. Note that $K_{\mathrm{pot}}(e)$ is the local contribution to the discrete Dirichlet energy $\mathcal{D}[u] = \tfrac{1}{2}\sum_{e \sim e'}[u(e')-u(e)]^2$, connecting the curvature indicator to a standard variational object.
\end{remark}

%==================================================================
\section{Compatibility and Boundary-Flux Identities}\label{sec:conservation}
%==================================================================

\begin{proposition}[Global source balance]\label{prop:balance}
The total source deviation vanishes:
\begin{equation}
    \sum_{e \in \Ecal} [S(e) - \bar{S}] = 0
\end{equation}
\end{proposition}

\begin{proof}
$\sum(S - \bar{S}) = \sum S - |\Ecal| \cdot \bar{S} = \sum S - \sum S = 0$.
\end{proof}

\begin{remark}
This is the compatibility identity for the centred source term in the graph Poisson equation, not a dynamical conservation law: it is forced by the definition of $\bar S$, independent of any field dynamics. Structurally, regions of positive source (excess content) are balanced by regions of negative source (deficit), and the field equation~\eqref{eq:field} redistributes this imbalance through the accessibility potential.
\end{remark}

\begin{proposition}[Discrete divergence theorem]\label{prop:flow}
Let $R \subset \Ecal$ be any (not necessarily connected) subset of the undirected Hasse graph $\Gcal$, and let
\[
    \partial^{\mathrm{out}} R \;=\; \{\,(e, e') \in \Ecal \times \Ecal : e \in R,\; e' \notin R,\; \{e,e'\} \in E_H\,\}
\]
be the \textit{oriented edge-boundary} of $R$: ordered pairs whose underlying undirected edge lies in $E_H$ and which point from inside $R$ to outside. For any function $u: \Ecal \to \mathbb{R}$ satisfying $\Delta_\Gcal u = S - \bar{S}$,
\begin{equation}
    \sum_{(e,e') \in \partial^{\mathrm{out}} R} [u(e') - u(e)] \;=\; \sum_{e \in R} [S(e) - \bar{S}].
\end{equation}
\end{proposition}

\begin{proof}
Sum the field equation over $e \in R$ and expand the Laplacian: $\sum_{e \in R} \sum_{e' \sim e} [u(e') - u(e)]$. For each undirected edge $\{e, e'\}$ with both endpoints in $R$, the two terms $[u(e') - u(e)]$ (from summing at $e$) and $[u(e) - u(e')]$ (from summing at $e'$) cancel. The remaining terms are exactly those with $e \in R$ and $e' \notin R$, i.e.\ $(e, e') \in \partial^{\mathrm{out}} R$. This is the discrete divergence theorem on an undirected graph.
\end{proof}

%==================================================================
\section{The Source--Curvature--Geodesic Bridge}\label{sec:bridge}
%==================================================================

The structural picture (to be read as a \emph{correspondence} in toy models, not a derived chain of theorems) is:
\begin{enumerate}
    \item \textbf{Source content} $S(e)$ measures a chosen local event-content deviation (density, signed overlap, or signed geodesic traffic).
    \item \textbf{The field equation} $\Delta_\Gcal u = S - \bar{S}$ determines the accessibility potential $u$ up to an additive constant (Proposition~\ref{prop:existence}).
    \item \textbf{Potential curvature} $K_{\mathrm{pot}}(e)$ is the squared discrete gradient of $u$ at $e$ (Definition~\ref{def:Kpot}).
    \item \textbf{Geodesics} in $\Gcal$ are independently defined as shortest paths (not modified by $u$); in the bottleneck toy model below, their concentration pattern is found to align with large $K_{\mathrm{pot}}$, but no geodesic equation is derived from $u$.
\end{enumerate}

Heuristically, this is the analogue of the Einstein programme
\[
    T_{\mu\nu} \;\xrightarrow{\text{field eq.}}\; G_{\mu\nu} \;\xrightarrow{\text{geodesic eq.}}\; \text{motion}
\]
replaced by the discrete pattern
\[
    S(e) \;\xrightarrow{\Delta_\Gcal u = S - \bar{S}}\; u(e),\, K_{\mathrm{pot}}(e) \;\xrightarrow{\text{observed correlation in toy models}}\; \text{geodesic concentration},
\]
with the final arrow being an observation about the examples computed below, not a theorem about general graphs.

%==================================================================
\section{Toy Model 1: Uniform Vacuum}\label{sec:vacuum}
%==================================================================

\begin{example}[Vacuum: symmetric $K_{5,5}$]\label{ex:vacuum}
The first-two-layer truncation of the dual-pentagon event set of Papers~XXIII--XXVI, with
\[
    L_1 = \{(0,4),(1,0),(2,1),(3,2),(4,3)\}, \qquad L_2 = \{(0,3),(1,4),(2,0),(3,1),(4,2)\},
\]
and Hasse graph $K_{5,5}$ (every $L_2$ event covers every $L_1$ event).

\textbf{Source.} Every vertex of $K_{5,5}$ has degree $5$; every ball $B_r(e)$ has the same size as every other ball of the same radius; so for any fixed $r$ and reference $V_{\mathrm{ref}} = \bar V$ the field $S_\rho(e)$ of Definition~\ref{def:rho} is identically zero. Similarly $S_\Omega(e) = \deg(e) - \overline{\deg} = 0$ and $S_J(e) = \beta(e) - \bar\beta = 0$ (vertex-transitivity). Thus $S - \bar S \equiv 0$ for all three source choices, and $\bar S = 0$.

\textbf{Field equation.} $\Delta_\Gcal u = 0$. Solution: $u \equiv \mathrm{const}$ (unique modulo constants).

\textbf{Curvature.} $K_{\mathrm{pot}}(e) = 0$ for every $e$: zero potential gradient everywhere.

\textbf{Geodesics.} By vertex-transitivity, every vertex carries the same geodesic betweenness; no concentration.

\textbf{Conclusion.} In this symmetric graph, zero source, constant potential, vanishing $K_{\mathrm{pot}}$, and uniform geodesic betweenness all coincide. The first three implications follow immediately from the displayed computation and Proposition~\ref{prop:existence}; the vanishing of betweenness non-uniformity follows independently from vertex-transitivity of $K_{5,5}$ (imported from Paper~XXVI).
\end{example}

%==================================================================
\section{Toy Model 2: Static Localised Source}\label{sec:static}
%==================================================================

\begin{example}[Localised source: bottleneck graph]\label{ex:bottleneck}
The three-layer bottleneck model from Paper~XXVI: 4 outer events in Layer~0, 1 bottleneck $b$ in Layer~1, 4 outer events in Layer~2.

\textbf{Source (overlap $S_\Omega$).} Degrees: $\deg(a_i) = 1$, $\deg(b) = 8$, $\deg(c_j) = 1$. Mean degree: $\overline{\deg} = 16/9$, so $\bar S_\Omega = 0$. Sources: $S_\Omega(b) = 8 - 16/9 = 56/9 \approx 6.22$; $S_\Omega(a_i) = S_\Omega(c_j) = 1 - 16/9 = -7/9 \approx -0.78$.

\textbf{Field equation.} $\Delta_\Gcal u = S_\Omega - \bar{S}_\Omega$. Since $\bar{S}_\Omega = 0$ (by construction of degree-deviation source), the equation is $\Delta_\Gcal u(e) = S_\Omega(e)$.

Solve on the bottleneck graph (9 nodes). By symmetry, $u(a_i) = u_a$ for all $i$ and $u(c_j) = u_c$ for all $j$. The equations are:
\begin{align}
    \Delta_\Gcal u(a_i) &= u(b) - u(a_i) = -7/9 \label{eq:bowtie-a}\\
    \Delta_\Gcal u(b) &= 4[u_a - u(b)] + 4[u_c - u(b)] = 56/9 \label{eq:bowtie-b}\\
    \Delta_\Gcal u(c_j) &= u(b) - u(c_j) = -7/9 \label{eq:bowtie-c}
\end{align}

From~\eqref{eq:bowtie-a} and~\eqref{eq:bowtie-c}: $u(b) - u_a = -7/9$ and $u(b) - u_c = -7/9$, so $u_a = u_c = u(b) + 7/9$.

Check~\eqref{eq:bowtie-b}: $4[(u(b) + 7/9) - u(b)] + 4[(u(b) + 7/9) - u(b)] = 4(7/9) + 4(7/9) = 56/9$. $\checkmark$

Setting $u(b) = 0$ (gauge choice): $u_a = u_c = 7/9$.

\textbf{Potential landscape.} The bottleneck sits at the \textit{lowest} potential ($u(b) = 0$); outer events sit higher ($u = 7/9$). The potential ``flows downhill'' toward the bottleneck.

\textbf{Curvature.} $K_{\mathrm{pot}}(b) = 4(7/9)^2 + 4(7/9)^2 = 8(49/81) = 392/81 \approx 4.84$. $K_{\mathrm{pot}}(a_i) = (7/9)^2 = 49/81 \approx 0.60$. The bottleneck has by far the highest potential-gradient curvature.

\textbf{Geodesics.} All shortest paths between non-bottleneck events pass through $b$ (Paper~XXVI). The potential gradient points toward $b$. Geodesic concentration and potential-flow direction are aligned.

\textbf{Conclusion.} This example exhibits a nontrivial source $S_\Omega$, a nonconstant Poisson solution $u$ with $K_{\mathrm{pot}}(b) \gg K_{\mathrm{pot}}(\text{leaves})$, and, independently (from Paper~XXVI), geodesic concentration through $b$. All three features coexist on this single toy graph; no general theorem linking the three is proved.
\end{example}

%==================================================================
\section{Toy Model 3: Perturbative Response}\label{sec:perturbation}
%==================================================================

\begin{example}[Perturbation of the flat model]\label{ex:perturb}
Start from the symmetric $K_{5,5}$ pentagon truncation of Example~\ref{ex:vacuum}. Add one extra event $f$ to $\Gcal$ with two new undirected edges: $f \sim (0,4)$ and $f \sim (1,0)$, attaching $f$ to two elements of $L_1$. The graph now has $11$ events and $27$ edges, is connected, and is no longer vertex-transitive.

\textbf{Symmetry classes of the perturbed graph.} The residual automorphism group of the perturbed $\Gcal$ consists of arbitrary permutations of $L_2$, the $\mathbb{Z}_2$ swap $(0,4) \leftrightarrow (1,0)$, and arbitrary permutations of the three unperturbed $L_1$ events $\{(2,1),(3,2),(4,3)\}$. Four orbits result:
\begin{itemize}
    \item $\mathcal{A} = \{(0,4),(1,0)\}$ (the two perturbed $L_1$ events, $|\mathcal{A}|=2$, $\deg = 6$);
    \item $\mathcal{B} = \{(2,1),(3,2),(4,3)\}$ (the unperturbed $L_1$ events, $|\mathcal{B}|=3$, $\deg = 5$);
    \item $\mathcal{C} = L_2$ (the five $L_2$ events, $|\mathcal{C}|=5$, $\deg = 5$);
    \item $\{f\}$ ($\deg = 2$).
\end{itemize}

\textbf{Source $S_\Omega$.} Mean degree $\overline{\deg} = (2\cdot 6 + 3\cdot 5 + 5\cdot 5 + 1\cdot 2)/11 = 54/11$, so $\bar S_\Omega = 0$. Sources by class:
\[
    S_\Omega(\mathcal{A}) = 6 - \tfrac{54}{11} = \tfrac{12}{11}, \;
    S_\Omega(\mathcal{B}) = S_\Omega(\mathcal{C}) = 5 - \tfrac{54}{11} = \tfrac{1}{11}, \;
    S_\Omega(f) = 2 - \tfrac{54}{11} = -\tfrac{32}{11}.
\]
Total: $2\cdot\tfrac{12}{11} + 8\cdot\tfrac{1}{11} + \bigl(-\tfrac{32}{11}\bigr) = 0$, consistent with Proposition~\ref{prop:balance}.

\textbf{Explicit Poisson solve.} On a perturbed $\Gcal$ respecting these four orbits, any solution of $\Delta_\Gcal u = S_\Omega$ is constant on each orbit. Let $u_\mathcal{A}, u_\mathcal{B}, u_\mathcal{C}, u_f$ denote the orbit values. Fix the gauge $u_f = 0$. Writing out the Laplacian at one representative of each orbit:
\begin{align}
    \Delta_\Gcal u \big|_{\mathcal{A}} &= 5 u_\mathcal{C} + u_f - 6 u_\mathcal{A} = \tfrac{12}{11}, \label{eq:pert-A}\\
    \Delta_\Gcal u \big|_{\mathcal{B}} &= 5 u_\mathcal{C} - 5 u_\mathcal{B} = \tfrac{1}{11}, \label{eq:pert-B}\\
    \Delta_\Gcal u \big|_{\mathcal{C}} &= 2 u_\mathcal{A} + 3 u_\mathcal{B} - 5 u_\mathcal{C} = \tfrac{1}{11}, \label{eq:pert-C}\\
    \Delta_\Gcal u \big|_{f} &= 2 u_\mathcal{A} - 2 u_f = -\tfrac{32}{11}. \label{eq:pert-f}
\end{align}
Equation~\eqref{eq:pert-f} gives $u_\mathcal{A} = -\tfrac{16}{11}$. Substituting in~\eqref{eq:pert-A} yields $u_\mathcal{C} = -\tfrac{84}{55}$, and~\eqref{eq:pert-B} then gives $u_\mathcal{B} = -\tfrac{17}{11}$. Equation~\eqref{eq:pert-C} is satisfied as a consistency check: $2\cdot\bigl(-\tfrac{16}{11}\bigr) + 3\cdot\bigl(-\tfrac{17}{11}\bigr) - 5\cdot\bigl(-\tfrac{84}{55}\bigr) = -\tfrac{32}{11} - \tfrac{51}{11} + \tfrac{84}{11} = \tfrac{1}{11}$. Summary: $u_f = 0$, $u_\mathcal{A} = -16/11 \approx -1.45$, $u_\mathcal{C} = -84/55 \approx -1.53$, $u_\mathcal{B} = -17/11 \approx -1.55$.

\textbf{Potential-gradient curvature.} Direct computation of $K_{\mathrm{pot}}(e) = \sum_{e' \sim e} [u(e') - u(e)]^2$ at one representative of each orbit:
\[
    K_{\mathrm{pot}}(f) = 2 (u_\mathcal{A} - u_f)^2 = \tfrac{512}{121} \approx 4.23,
\]
\[
    K_{\mathrm{pot}}(\mathcal{A}) = 5 (u_\mathcal{C} - u_\mathcal{A})^2 + (u_f - u_\mathcal{A})^2 = \tfrac{1296}{605} \approx 2.14,
\]
\[
    K_{\mathrm{pot}}(\mathcal{B}) = 5 (u_\mathcal{C} - u_\mathcal{B})^2 = \tfrac{1}{605} \approx 0.0017,
\]
\[
    K_{\mathrm{pot}}(\mathcal{C}) = 2 (u_\mathcal{A} - u_\mathcal{C})^2 + 3 (u_\mathcal{B} - u_\mathcal{C})^2 = \tfrac{7}{605} \approx 0.0116.
\]
The perturbation event $f$ has the largest potential-gradient curvature; the two attached events $\mathcal{A}$ have the next largest. The unperturbed classes $\mathcal{B}, \mathcal{C}$ have $K_{\mathrm{pot}}$ three orders of magnitude smaller.

\textbf{Geodesics.} The new edges open a length-$2$ shortest path $f \to (0,4) \to L_2$ and a length-$3$ shortest path $f \to (0,4) \to L_2 \to \mathcal{B}$-event; the existing $K_{5,5}$ geodesics between $L_1$ and $L_2$ are unchanged. A full betweenness recomputation on the $11$-vertex perturbed graph is not given here.

\textbf{Conclusion.} Adding a single event to $K_{5,5}$ produces a closed-form Poisson solution with nonzero source at four orbits, a potential landscape with a localised well at $\mathcal{B}$ and a pronounced gradient at $f$, and $K_{\mathrm{pot}}$ concentrated at the perturbation event. In this example the source-to-potential-to-$K_{\mathrm{pot}}$ arithmetic is fully checked; the geodesic-level claims are restricted to the existence of the new shortest paths exhibited above, not to a betweenness redistribution theorem.
\end{example}

%==================================================================
\section{Relation to Einstein-Style Dynamics}\label{sec:einstein}
%==================================================================

\begin{table}[ht]
\centering
\caption{Structural comparison: VFD field dynamics vs.\ Einstein gravity.}
\label{tab:einstein}
\begin{tabular}{@{}lll@{}}
\toprule
Feature & VFD (this paper) & General relativity \\
\midrule
Source & Local event-content deviation $S(e)$ & Stress-energy $T_{\mu\nu}$ \\
Field variable & Accessibility potential $u(e)$ & Metric $g_{\mu\nu}$ \\
Field equation & $\Delta_\Gcal u = S - \bar{S}$ & $G_{\mu\nu} = 8\pi G\, T_{\mu\nu}$ \\
Curvature & Potential gradient $K_{\mathrm{pot}}$ & Riemann tensor $R^\alpha{}_{\beta\gamma\delta}$ \\
Compatibility / conservation & $\sum_e [S(e) - \bar{S}] = 0$ (algebraic) & $\nabla_\mu T^{\mu\nu} = 0$ (dynamical) \\
Geodesic response & Shortest-path concentration & Geodesic equation \\
Vacuum solution & $u = \mathrm{const}$, flat & Minkowski spacetime \\
\bottomrule
\end{tabular}
\end{table}

\begin{remark}[Epistemic status]
The structural parallels in Table~\ref{tab:einstein} are exhibited, not derived from the Einstein equations. The VFD field equation $\Delta_\Gcal u = S - \bar{S}$ is a discrete Poisson equation on a graph, a standard mathematical object~\citep{Chung1997}. Its physical interpretation as a gravity analogue is structural: it shows that the event-order framework supports a source-to-potential construction and, in the bowtie toy model, a correlation between large potential gradients and concentrated geodesic traffic. No motion law is derived from $u$, and no continuum limit is claimed. Whether this discrete graph-Poisson model converges to any continuum Einstein-type limit in an appropriate refinement is an open question of significant interest.
\end{remark}

%==================================================================
\section{Discussion and Limitations}\label{sec:discussion}
%==================================================================

\subsection{What is established}

\begin{itemize}
    \item Three source candidates constructed: density, overlap, geodesic traffic (Definitions~\ref{def:rho}--\ref{def:traffic}).
    \item Accessibility potential defined via graph Laplacian (Definition~\ref{def:potential}).
    \item Discrete field equation $\Delta_\Gcal u = S - \bar{S}$: existence and uniqueness (up to constant) proved (Proposition~\ref{prop:existence}).
    \item Global balance law: total source deviation vanishes (Proposition~\ref{prop:balance}).
    \item Integrated boundary-flux identity for solutions of the Poisson equation: discrete divergence theorem (Proposition~\ref{prop:flow}).
    \item Potential-gradient curvature $K_{\mathrm{pot}}$ defined and computed (Definition~\ref{def:Kpot}).
    \item Vacuum model: flat, zero source, constant potential, uniform geodesics.
    \item Bottleneck model: localised source, nontrivial potential, geodesic concentration.
    \item Perturbation model: single-event addition creates measurable geometric response.
    \item Partial source-to-potential correspondence exhibited explicitly, with geodesic alignment verified only in the bowtie example.
\end{itemize}

\subsection{What is not established}

\begin{itemize}
    \item No Einstein equations. The field equation is a graph Poisson equation, not the Einstein field equations.
    \item No stress-energy tensor. The source $S(e)$ is a scalar graph quantity, not a tensor.
    \item No continuum limit. Whether $\Delta_\Gcal u = S - \bar{S}$ converges to $G_{\mu\nu} = 8\pi G T_{\mu\nu}$ is open.
    \item No time evolution. The field equation is solved statically; dynamical evolution of source and geometry is not addressed.
    \item No coupling constant. The proportionality between source and geometric response is set to 1; a Newton-constant analogue is not derived.
    \item The choice of source field is not uniquely determined by the framework; density, overlap, and geodesic-traffic constructions are all viable candidates.
\end{itemize}

\subsection{Programme position}

\begin{itemize}
    \item Paper~XXIII: event ordering (derived); the no-canonical-clock conclusion is conditional on incomparable events, a hypothesis verified only in the pentagon model.
    \item Paper~XXIV: observer kinematics; observer definition unconditional, simultaneity conditional on incomparable events, endpoint-interval (time-dilation) theorem conditional on unequal weight sums along a shared linear extension.
    \item Paper~XXV: separation geometry; global metric on the event set requires connectedness of the Hasse graph (verified for the pentagon truncation, open for the full 600-cell set).
    \item Paper~XXVI: three candidate graph-curvature indicators; no comparison theorem among them is proved.
    \item Paper~XXVII (this paper): discrete graph-Poisson analogue and a potential-gradient curvature indicator, exhibited on three toy graphs; no Einstein equation, no stress-energy tensor, no continuum limit, and no geodesic equation derived from $u$.
\end{itemize}

%==================================================================
\section{Conclusion}\label{sec:conclusion}
%==================================================================

The event-order framework is equipped with a discrete graph-Poisson source-to-potential model: the graph Laplacian maps a signed source field $S - \bar S$ (from the chosen source family) to an accessibility potential $u$ on the Hasse graph, determined up to an additive constant. In the bowtie toy model, the associated potential-gradient curvature $K_{\mathrm{pot}}$ is maximal at the same bottleneck $b$ where Paper~XXVI found maximal $K_{\mathrm{branch}}$ and $K_{\mathrm{geo}}$; the $K_{\mathrm{pot}}$--$K_{\mathrm{branch}}$ alignment is a same-degree-data coincidence (both built from $\deg - \overline{\deg}$), while the $K_{\mathrm{pot}}$--$K_{\mathrm{geo}}$ co-location is not. In the perturbation toy model, closed-form $K_{\mathrm{pot}}$ values are computed by symmetry reduction, but a geodesic-betweenness recomputation on the perturbed graph is not provided. The field equation is introduced as a model definition, not deduced from the earlier VFD papers; no Einstein equation, no stress-energy tensor, no geodesic equation derived from $u$, and no continuum limit is established.

At the level proved so far, the GR-leg programme now has a discrete scaffold for each of the following structural themes, subject to the standing conditionals inherited from the earlier papers:
\begin{itemize}
    \item \textbf{Time}: event ordering from phase-overlap admissibility (XXIII; no-canonical-clock conditional on incomparable events, verified in the pentagon model).
    \item \textbf{Observers}: frame-dependent kinematics from linear extensions (XXIV; simultaneity conditional on incomparable events, endpoint-interval theorem conditional on unequal weight sums along a shared linear extension).
    \item \textbf{Metric}: separation constructions including a transition-cost metric on connected components of the Hasse graph (XXV).
    \item \textbf{Curvature}: three candidate graph indicators on the Hasse graph, no comparison theorem (XXVI).
    \item \textbf{Source-field analogue}: the discrete Poisson equation $\Delta_\Gcal u = S - \bar S$ and the associated potential-gradient indicator $K_{\mathrm{pot}}$ (this paper).
\end{itemize}

Whether any of this discrete scaffold converges to the continuum Einstein dynamics in an appropriate refinement limit is the defining open question for the next stage of the programme.

%==================================================================
\bibliographystyle{plainnat}
\bibliography{references}

\end{document}
